%!TEX root = ../hutbthesis_main.tex

\clearpage
\phantomsection
\addcontentsline{toc}{chapter}{参考文献}
\begingroup
\songti
\zihao{-4}
\setlength{\parskip}{0pt}
\linespread{1.0}\selectfont
\sloppy

\begin{thebibliography}{99}
	\vspace{-1.5em}
	\setlength{\itemsep}{0pt}
	\setlength{\parsep}{0pt}
	\setlength{\parskip}{0pt}
\bibitem{ref1} 陈夏非,焦翊洋,姜雪,等.基于自适应模型预测控制的驾驶模拟器运动提示算法设计与验证[J/OL].机械工程学报,1-10[2026-04-28].https://link.cnki.net/urlid/11.2187.th.20260225.1756.008.
\bibitem{ref2} 孔令智,田毕江,熊昌安,等.汽车驾驶模拟器的应用研究综述[J].公路与汽运,2025,41(04):1-9.DOI:10.20035/j.issn.1671-2668.2025.04.001.
\bibitem{ref3} 王俊达,卿杜政.柔性可扩展装备体系对抗仿真建模框架研究[J].现代防御技术,2020,48(04):122-131.
\bibitem{ref4}
Liuzzo M S, Carmignani N, Carver L, et al. Commissioning simulations tools based on python Accelerator Toolbox[J]. Journal of Physics: Conference Series, 2024, 2687(3): DOI:10.1088/1742-6596/2687/3/032001.
\bibitem{ref5} 杨博,张能,李善平,等.智能代码补全研究综述[J].软件学报,2020,31(05):1435-1453.DOI:10.13328/j.cnki.jos.005966.
\bibitem{ref6} 许鹏宇,况博裕,苏铓,等.基于大语言模型的自动代码修复综述[J].计算机研究与发展,2025,62(08):2040-2057.
\bibitem{ref7} 郭家顺.基于LLM的智能代码补全系统设计与性能评估[J].软件,2025,46(10):35-37.
\bibitem{ref8} 李溪子.基于机器学习算法的智能代码补全工具开发与实现[J].软件,2025,46(02):98-100.
\bibitem{ref9} 邹佰翰,汪莹,彭鑫,等.重新审视代码补全中的检索增强策略[J].软件学报,2025,36(06):2747-2773.DOI:10.13328/j.cnki.jos.007226.
\bibitem{ref10} 陈伟,何成万,余秋惠,等.基于CNN和自注意力神经网络的代码补全方法[J].计算机工程与设计,2025,46(10):2919-2926.DOI:10.16208/j.issn1000-7024.2025.10.024.
\bibitem{ref11} 杨泽洲.基于语言模型的代码构建辅助技术研究[D].哈尔滨工业大学,2025.DOI:10.27061/d.cnki.ghgdu.2025.002727.
\bibitem{ref12} 许柏炎,蔡瑞初,梁智豪.一种用于代码注释自动生成的语法辅助复制机制[J].计算机工程,2021,47(04):92-99.DOI:10.19678/j.issn.1000-3428.0057290.
\bibitem{ref13} 王莹莹.基于预训练模型的目标描述代码自动生成研究[D].河北师范大学,2024.DOI:10.27110/d.cnki.ghsfu.2024.000468.
\bibitem{ref14} Hostnik M ,Šikonja R M .Retrieval-augmented code completion for local projects using large language models[J].Expert Systems With Applications,2025,292128596-128596.DOI:10.1016/J.ESWA.2025.128596.
\bibitem{ref15} Husein A R ,Aburajouh H ,Catal C .Large language models for code completion: A systematic literature review[J].Computer Standards \& Interfaces,2025,92103917-103917.DOI:10.1016/J.CSI.2024.103917.
\bibitem{ref16}
Wannita T, Chakkrit T, Yuan-Fang L. Syntax-aware on-the-fly code completion[J]. Information and Software Technology, 2024, 165: DOI:10.1016/J.INFSOF.2023.107336.
\bibitem{ref17} Abufarha S ,Marouf A A ,Rokne G J , et al.Mitigating Prompt Dependency in Large Language Models: A Retrieval-Augmented Framework for Intelligent Code Assistance[J].Software,2026,5(1):4-4.DOI:10.3390/SOFTWARE5010004.
\bibitem{ref18} Chang S ,Geng C ,Huang H , et al.CodeSpeak: Improving smart contract vulnerability detection via LLM-assisted code analysis[J].The Journal of Systems \& Software,2026,231112635-112635.DOI:10.1016/J.JSS.2025.112635.
\bibitem{ref19} 刘昌,周长鑫.BIM模型编码插件的开发与应用[J].工程技术研究,2023,8(19):124-126.DOI:10.19537/j.cnki.2096-2789.2023.19.041.
\bibitem{ref20} 茆凯成,康彦.基于深度学习技术的智能对话插件开发研究——以心理健康领域为例[J].中国新通信,2024,26(24):127-129.
\end{thebibliography}
\endgroup
